\documentclass[11pt,a4paper]{article}

\usepackage[utf8]{inputenc}
\usepackage[T1]{fontenc}
\usepackage[margin=1in]{geometry}
\usepackage{amsmath,amssymb}
\usepackage{booktabs}
\usepackage{graphicx}
\usepackage{hyperref}
\usepackage{microtype}

\newcommand{\TVSumVideos}{50}
\newcommand{\OracleAnnotators}{20}
\newcommand{\DSNetUniqueTestVideos}{32}
\newcommand{\DSNetSplitCount}{5}
\newcommand{\AblationArmCount}{6}
\newcommand{\BarycenterDimension}{one-dimensional}
\newcommand{\BootstrapConfidenceLevel}{95}
\newcommand{\OracleDirection}{is lower than the arithmetic mean on F1 and both rank-correlation averages}
\newcommand{\AblationUniformFinding}{None of the A--F ablation arms has a paired test against uniform below the conventional significance threshold.}
\newcommand{\BudgetPercent}{15}
\newcommand{\BootstrapReplicates}{10000}
\newcommand{\SimulationAnnotators}{20}
\newcommand{\SimulationShiftWidths}{0, 1, 2, 4, 8, 16, 32}
\newcommand{\SimulationNoisePercent}{15.0}
\newcommand{\SimulationFirstPositiveCrossover}{not observed}
\newcommand{\BarycenterW2Label}{characterization only; barycenter minimizes W2 by construction}
\newcommand{\ExploratorySigma}{4}
\newcommand{\BilstmUniformDifference}{2.492}
\newcommand{\BilstmUniformTTestP}{0.0493}
\newcommand{\BilstmUniformWilcoxonP}{0.0683}


\title{When Does Wasserstein Consensus Help? Annotator Disagreement in Video Summarization on TVSum}
\author{Research Team\\Indian Institute of Information Technology, Kottayam}
\date{}

\begin{document}
\maketitle

\begin{abstract}
We study whether a \BarycenterDimension{} Wasserstein consensus target helps
video-conditioned summarization when human annotators disagree. The analysis
separates leave-one-annotator-out oracle evaluations from learned models and
from a labelled jitter simulation. The arithmetic mean and quantile-averaged
barycenter are compared on the real TVSum annotations and preprocessed visual
features. The barycenter minimizes the chosen Wasserstein objective by
construction, but oracle summaries and learned-model results do not show a
general downstream advantage. Mechanism analyses are associative, and the
jitter experiment is explicitly simulated rather than additional real-data
evidence. All numerical tables are generated from saved run JSON files.
\end{abstract}

\section{Introduction}
Video summarization requires translating frame-level importance judgments
into a compact set of shots. With multiple annotators, a target distribution
can be formed by averaging their normalized ratings in the original pick
index, or by computing a Wasserstein barycenter. These targets optimize
different notions of consensus: an arithmetic mean minimizes squared
Euclidean distance, while a \BarycenterDimension{} quantile-averaged barycenter
minimizes average squared Wasserstein distance on its support.

This work asks when that transport-based consensus helps downstream
summarization, rather than whether it solves its defining transport
objective. We use the real TVSum videos, their annotator ratings, and the
saved video-conditioned benchmark outputs. We distinguish oracle target
evaluation, learned prediction, descriptive mechanism analysis, and a
controlled but labelled simulation. Results are generated from saved JSON
artifacts by \texttt{build\_tables.py}; the corresponding input keys and
checks are recorded in \texttt{tables/table\_sources.json}.

\section{Related Work}
\textbf{TODO:} Add a verified review of video summarization and multi-annotator
consensus methods. We cite the original TVSum dataset paper
\cite{song2015tvsum}; transport and diffusion references should be added only
after bibliographic details and their relevance are independently verified.

\section{Consensus Targets and the Leave-One-Annotator-Out Protocol}
For each video, raw annotator ratings are aligned to the preprocessed feature
pick indices and normalized to probability distributions. In a
leave-one-annotator-out (LOAO) trial, one annotator is held out and the
remaining distributions define the target. The arithmetic target is their
pointwise mean. The barycenter is the \BarycenterDimension{} quantile-averaged
Wasserstein barycenter on pick-index support, with interpolated mass
aggregated back to the original picks. The held-out annotator's original-frame
summary is used for evaluation. The saved oracle protocol averages over
\OracleAnnotators{} leave-one-annotator-out trials per video.

Predicted pick scores are expanded to frames, averaged within the supplied
shot boundaries, and converted to a summary by the canonical knapsack
procedure at a frame budget of \BudgetPercent\%. Oracle results average each metric over held-out
annotators within a video, then pair the resulting video-level values. The
learned experiment follows the supplied DSNet splits, with validation
restricted to the training portion. The split file has incomplete unique
test-video coverage, so reported learned-model tests use the saved common
unique-video pairing and are not a disjoint all-video cross-validation
estimate.

\section{Oracle Results}
The saved LOAO oracle results compare the arithmetic mean with the
quantile-averaged barycenter on F1 and rank correlation. Paired tests are
performed over videos, and the reported adjusted values use the saved Holm
correction where specified. On the saved oracle metrics, the barycenter
\OracleDirection{}.
\begin{table}[ht]
\centering
\caption{Leave-one-annotator-out real-TVSum oracle metrics. Mean and barycenter estimates are averages over videos; paired tests compare per-video means.}
\label{tab:oracle}
\resizebox{\textwidth}{!}{%
\begin{tabular}{lrrrrrr}
\toprule
Metric & Mean & Barycenter & Difference & Paired $t$ $p$ (raw/Holm) & Wilcoxon $p$ (raw/Holm) \\
\midrule
F1 (\%) & 62.621 & 60.681 & -1.940 & 0.00113/0.00604 & 0.000475/0.00244 \\
Spearman $\rho$ & 0.395 & 0.244 & -0.151 & 1.13e-22/9.06e-22 & 1.78e-15/1.42e-14 \\
Kendall $\tau$ & 0.309 & 0.192 & -0.117 & 2.95e-22/2.36e-21 & 1.78e-15/1.42e-14 \\
\bottomrule
\end{tabular}%
}
\vspace{1mm}
{\footnotesize Paired test columns show raw/Holm-adjusted p-values. F1 uses the saved oracle paired comparison; rank metrics use Checkpoint secondary paired tests.}
\end{table}

\begin{table}[ht]
\centering
\caption{Secondary LOAO metrics for mean and barycenter targets. Paired differences are barycenter minus mean; p-values show raw/Holm adjustment.}
\label{tab:secondary}
\resizebox{\textwidth}{!}{%
\begin{tabular}{lrrrrr}
\toprule
Metric & Mean & Barycenter & Difference & Paired $t$ $p$ (raw/Holm) & Wilcoxon $p$ (raw/Holm) \\
\midrule
Spearman $\rho$ & 0.395 & 0.244 & -0.151 & 1.13e-22/9.06e-22 & 1.78e-15/1.42e-14 \\
Kendall $\tau$ & 0.309 & 0.192 & -0.117 & 2.95e-22/2.36e-21 & 1.78e-15/1.42e-14 \\
Tolerance F1, $k=2$ & 69.262 & 67.955 & -1.306 & 0.00638/0.0319 & 0.00182/0.0091 \\
Tolerance F1, $k=5$ & 76.589 & 76.218 & -0.371 & 0.384/1 & 0.137/0.742 \\
Tolerance F1, $k=10$ & 81.905 & 81.830 & -0.075 & 0.856/1 & 1/1 \\
\bottomrule
\end{tabular}%
}
\end{table}


The barycenter minimizes average \BarycenterDimension{} squared Wasserstein distance
by construction. Its W2 comparison against held-out normalized annotator
distributions is therefore reported as characterization only, not as
evidence of improved summary quality.
\begin{table}[ht]
\centering
\caption{Wasserstein-distance characterization only; the barycenter minimizes W2 by construction.}
\label{tab:w2}
\resizebox{\textwidth}{!}{%
\begin{tabular}{rrrrrr}
\toprule
Mean-target W2 & Barycenter W2 & Difference & Paired $t$ $p$ raw/Holm & Wilcoxon $p$ raw/Holm & $n$ \\
\midrule
16.097 & 15.984 & -0.113 & 1.22e-09/1.22e-09 & 1.95e-13/1.95e-13 & 50 \\
\bottomrule
\end{tabular}%
}
\end{table}


\section{Learned-Model Results}
The saved ablation evaluates \AblationArmCount{} arms, labelled A--F, from the
temporal-convolution run. The additional BiLSTM comparison uses the saved
mean-target, cross-entropy arm. For those ablation arms, both paired tests against the
uniform summary are reported; neither test family is silently substituted
for the other. \AblationUniformFinding{} For the BiLSTM, the saved-data paired
comparison is mixed: the paired $t$-test gives $p=\BilstmUniformTTestP{}$,
while Wilcoxon gives $p=\BilstmUniformWilcoxonP{}$. Thus these results do
not support a consistent claim that learned models exceed uniform
summarization. The BiLSTM-versus-uniform comparison is reconstructed from the
saved backbone and uniform per-video records using the same unique-video
aggregation and paired-test procedure.
\begin{table}[ht]
\centering
\caption{Learned-model metrics generated from the saved real-TVSum JSON. F1 is a percentage; rank correlations are unitless.}
\label{tab:learned-metrics}
\begin{tabular}{lrrrr}
\toprule
Model & F1 mean & Seed SD & Spearman $\rho$ & Kendall $\tau$ \\
\midrule
A: Mean target + MSE & 60.766 & 2.513 & 0.1479 & 0.1126 \\
B: Mean target + cross-entropy & 61.216 & 2.483 & 0.1682 & 0.1281 \\
C: Mean target + Cramer loss & 60.283 & 2.949 & 0.1352 & 0.1028 \\
D: Barycenter target + MSE & 60.371 & 2.694 & 0.1157 & 0.0881 \\
E: Barycenter target + Cramer loss & 60.600 & 2.923 & 0.1436 & 0.1093 \\
F: Validation-selected blend + Cramer loss & 60.549 & 2.866 & 0.1436 & 0.1092 \\
BiLSTM, mean target + cross-entropy & 61.860 & 2.487 & 0.1992 & 0.1520 \\
\bottomrule
\end{tabular}
\end{table}

\begin{table}[ht]
\centering
\caption{Paired learned-model versus uniform-summary F1 comparisons over the saved unique test videos. Raw paired-test p-values are shown.}
\label{tab:learned-uniform}
\begin{tabular}{lrrrr}
\toprule
Learned model & $\Delta$ F1 & $n$ & Paired $t$ $p$ & Wilcoxon $p$ \\
\midrule
A: Mean target + MSE & 1.283 & 32 & 0.253 & 0.345 \\
B: Mean target + cross-entropy & 2.015 & 32 & 0.0814 & 0.0741 \\
C: Mean target + Cramer loss & 0.864 & 32 & 0.474 & 0.733 \\
D: Barycenter target + MSE & 0.749 & 32 & 0.521 & 0.775 \\
E: Barycenter target + Cramer loss & 1.141 & 32 & 0.337 & 0.443 \\
F: Validation-selected blend + Cramer loss & 1.242 & 32 & 0.297 & 0.421 \\
BiLSTM, mean target + cross-entropy & 2.492 & 32 & 0.0493 & 0.0683 \\
\bottomrule
\end{tabular}
\end{table}

\begin{table}[ht]
\centering
\caption{Saved paired BiLSTM-minus-temporal-convolution comparison for arm B.}
\label{tab:bilstm}
\begin{tabular}{lrrrr}
\toprule
Metric & Difference & $n$ & Paired $t$ $p$ & Wilcoxon $p$ \\
\midrule
F1 (points) & 0.477 & 32 & 0.153 & 0.507 \\
Spearman $\rho$ & 0.031 & 32 & 3.44e-05 & 3.54e-05 \\
Kendall $\tau$ & 0.024 & 32 & 2.68e-05 & 2.54e-05 \\
\bottomrule
\end{tabular}
\end{table}


\section{Disagreement and Mechanism Analysis}
We report associations between the saved per-video barycenter-minus-mean F1
gap and only the predeclared descriptors: interior peaks in the normalized
mean consensus, saved mean pairwise annotator Spearman correlation,
dominant-peak location dispersion, and mass-weighted overlap of canonical
selections. Correlation intervals are percentile bootstrap intervals over
videos. They are associations, not causal effects. Qualitative examples are
summarized from their saved per-video records.
\begin{table}[ht]
\centering
\caption{Associations between predeclared per-video descriptors and barycenter-minus-mean F1. All entries are descriptive associations.}
\label{tab:mechanism}
\begin{tabular}{lrrr}
\toprule
Descriptor & Spearman $\rho$ & Bootstrap \BootstrapConfidenceLevel\% CI & Videos \\
\midrule
Mean-curve peak count & -0.310 & [-0.572, -0.017] & 50 \\
Annotator pairwise Spearman & -0.289 & [-0.511, -0.028] & 50 \\
Peak-location SD (picks) & -0.210 & [-0.461, 0.072] & 50 \\
Weighted selected-mass overlap & 0.293 & [0.009, 0.542] & 50 \\
\bottomrule
\end{tabular}
\end{table}

\begin{table}[ht]
\centering
\caption{Saved qualitative case-study measurements. These selected videos are descriptive examples.}
\label{tab:mechanism-cases}
\begin{tabular}{lrrrr}
\toprule
Video & Mean F1 & Barycenter F1 & Difference & Shared pick count \\
\midrule
uGu\_10sucQo & 41.838 & 51.800 & 9.962 & 41 \\
i3wAGJaaktw & 74.410 & 64.329 & -10.081 & 38 \\
\bottomrule
\end{tabular}
\end{table}


An exploratory smoothing analysis uses the saved high-disagreement subset.
The $\sigma=\ExploratorySigma{}$ setting is reported as exploratory and was not used to tune or
select a method.
\begin{table}[ht]
\centering
\caption{Exploratory high-disagreement smoothing gains; no smoothing value was used for tuning or selection.}
\label{tab:exploratory-smoothing}
\begin{tabular}{lrrrr}
\toprule
Variant & Mean F1 gain & Bootstrap \BootstrapConfidenceLevel\% CI & Improved videos / total & Videos \\
\midrule
Mean smoothing, $\sigma=1$ & 0.118 & [-0.534, 0.708] & 8 / 17 & 17 \\
Mean smoothing, $\sigma=2$ & 0.766 & [-0.241, 1.740] & 11 / 17 & 17 \\
Mean smoothing, $\sigma=4$ & 1.498 & [0.657, 2.348] & 12 / 17 & 17 \\
\bottomrule
\end{tabular}
\end{table}


\section{Controlled Jitter Simulation}
The simulation takes each real video's saved-source mean consensus curve,
generates simulated annotator curves using the declared random shifts and
noise, and compares LOAO consensus summaries using the same canonical
summary procedure. The saved simulation uses \SimulationAnnotators{} simulated
annotators per video and shift width, shift half-widths of
\SimulationShiftWidths{} picks, and noise with standard deviation
\SimulationNoisePercent\% of the source curve maximum. The percentile
intervals resample videos \BootstrapReplicates{} times. Every output is
labelled simulation and excluded from the real-data results. A positive mean
barycenter-gain crossover was \SimulationFirstPositiveCrossover{} in the
declared range; this controlled result is not evidence about the actual
annotators.
\begin{table}[ht]
\centering
\caption{SIMULATION ONLY: labelled jitter experiment on real-video source curves; these results are not real annotator measurements.}
\label{tab:simulation}
\resizebox{\textwidth}{!}{%
\begin{tabular}{rrrrrr}
\toprule
Shift half-width (picks) & Mean F1 & Barycenter F1 & Gain & Bootstrap \BootstrapConfidenceLevel\% CI & Improved videos \\
\midrule
0 & 86.553 & 85.651 & -0.902 & [-1.686, -0.121] & 10 / 50 \\
1 & 85.747 & 84.717 & -1.030 & [-1.785, -0.363] & 7 / 50 \\
2 & 84.663 & 83.941 & -0.722 & [-1.359, -0.094] & 11 / 50 \\
4 & 83.211 & 82.867 & -0.344 & [-0.994, 0.222] & 12 / 50 \\
8 & 80.342 & 79.054 & -1.288 & [-2.487, -0.404] & 10 / 50 \\
16 & 76.445 & 75.691 & -0.754 & [-1.956, 0.400] & 19 / 50 \\
32 & 73.340 & 72.410 & -0.930 & [-1.848, -0.052] & 13 / 50 \\
\bottomrule
\end{tabular}%
}
\end{table}


\section{Limitations}
TVSum contains \TVSumVideos{} videos and the saved DSNet experiment covers
\DSNetUniqueTestVideos{} unique test videos under \DSNetSplitCount{} supplied
splits. Its overlapping test sets and incomplete coverage limit
generalization. The barycenter studied here is the
\BarycenterDimension{}
quantile-averaged version on pick-index support; other discretizations or
consensus objectives are outside this analysis. Oracle and learned results
use distinct protocols and should not be conflated. Mechanism associations
do not identify causal explanations, and the controlled jitter experiment is
labelled simulation. The exploratory smoothing result was not used to tune
the benchmark. No claim of improved summarization, state of the art, or
diffusion-model performance is made.

\section{Conclusion}
On this benchmark, minimizing Wasserstein distance does not by itself imply
better F1 or rank agreement. The saved oracle and learned results, along with
the explicitly separated mechanism and simulation analyses, motivate a
conditional rather than universal account of when consensus choices may
matter. Any broader claim requires independently validated data and
evaluation.

\bibliographystyle{plain}
\bibliography{references}

\end{document}
